\section*{Розділ I. Загальні положення}
\addcontentsline{toc}{section}{Розділ I. Загальні положення}

\subsection*{1.1. Правова основа та мета Регламенту}
\addcontentsline{toc}{subsection}{1.1. Правова основа та мета Регламенту}
    1.1.1. Регламент Конференції студентів Університету (далі -- Регламент КСУ) розроблено відповідно до Конституції України, Закону України "Про вищу освіту", Статуту Київського авіаційного інституту та Положення про органи студентського самоврядування Київського авіаційного інституту.

    1.1.2. Цей Регламент визначає процедурні засади організації та проведення засідань Конференції студентів Університету (КСУ), порядок прийняття рішень, а також регулює інші процедурні питання діяльності КСУ.

    1.1.3. Регламент КСУ є обов'язковим для виконання всіма делегатами КСУ, Спікером КСУ, запрошеними особами та іншими учасниками засідань КСУ.

\subsection*{1.2. Основні терміни та визначення}
\addcontentsline{toc}{subsection}{1.2. Основні терміни та визначення}
    1.2.1. У цьому Регламенті терміни вживаються у значеннях, визначених Положенням про органи студентського самоврядування Київського авіаційного інституту та глосарії цього Регламенту.

    1.2.2. Для цілей цього Регламенту використовуються такі додаткові процедурні терміни:

        \begin{enumerate}[label=\alph*)]
            \item \textbf{Кворум КСУ} -- мінімальна кількість делегатів КСУ (не менше половини від загального складу), необхідна для правомочності засідання КСУ;
            \item \textbf{Порядок денний} -- перелік питань, що підлягають розгляду на засіданні КСУ, затверджений у встановленому цим Регламентом порядку;
            \item \textbf{Робочий орган КСУ} -- тимчасова комісія або інший допоміжний орган, створений КСУ для виконання окремих завдань або підготовки питань.
        \end{enumerate}

\subsection*{1.3. Принципи роботи КСУ}
\addcontentsline{toc}{subsection}{1.3. Принципи роботи КСУ}
    1.3.1. Діяльність КСУ здійснюється на принципах колегіальності, виборності делегатів, рівності прав делегатів, гласності та відкритості, визначених пунктом 3.1.3 Положення про органи студентського самоврядування Київського авіаційного інституту.

    1.3.2. Процедурні аспекти реалізації зазначених принципів визначаються цим Регламентом.

\subsection*{1.4. Дія Регламенту у часі}
\addcontentsline{toc}{subsection}{1.4. Дія Регламенту у часі}
    1.4.1. Цей Регламент затверджується Конференцією студентів Університету (КСУ) простою більшістю голосів від загального складу делегатів КСУ.

    1.4.2. Регламент КСУ набуває чинності з дня його офіційного оприлюднення на інформаційних ресурсах органів студентського самоврядування Університету після його затвердження КСУ.

    1.4.3. Зміни та доповнення до цього Регламенту вносяться у порядку, визначеному Розділом VIII цього Регламенту.